\begin{afterword}
\summaryhead

The post answers Paul Gowder, who had argued that the author's preference for one person's torture over a googolplex dust specks depends on utilitarian intuitions that the author has not argued for. The author accepts that all moral reasoning starts from intuitions, since there is nothing else to argue from. The task is to build on the intuitions that survive reflection; keeping every particular one leads to circular preferences.

The author separates intuitions about where to go from intuitions about how to get there, and holds that the second kind are badly flawed. People give more to save one child than to save a group of eight children, and pay no more to save ten times as many lives. The explanation is that the brain does not multiply. If saving 50,000 lives is better than saving 25,000, and a life's value does not depend on how many others exist, values add up. Sacred values and absolute rules are explained as signals of commitment, not as evidence for two tiers of value; a lexically lower tier would decide nothing. The value of a single event may be complex, but when lives are at stake one should ``shut up and multiply''.

\responsehead

On method, the post gives a good answer to Gowder. It grants that intuitions are all there is and asks which of them survive reflection. That is close to reflective equilibrium, the dominant method in moral philosophy, and it moves the dispute to the right question.

Its evidence that people fail to count is well chosen and fairly reported. The Slovic quotation matches its source. The study behind it found what the post says it found, and explained it, as the post does, by the emotion one identified child arouses. A second paper by the same authors supports the post's reading: when people had to choose between the one child and the group, 43 of 62 chose the group.

That evidence, though, concerns quantities of the same good, lives or children. The post's arguments for aggregation (five efforts against one, ``Three lives are one and one and one'') are of the same kind. Many contractualists, who follow Scanlon in refusing to add up complaints, still hold that a rescuer should save five people rather than one. What they deny is that small harms to many can outweigh a grave harm to one; Parfit wrote that we might save one person from a hundred days of pain rather than ``any number'' of others from ten. So this evidence, and these arguments, do not address the case Gowder raised.

The post's argument that does address it is the one against two tiers of value: a lexically lower tier would decide nothing, because some upper-tier consideration is always present. This is a real objection. Huemer later developed it against lexical priority in general, and the Stanford Encyclopedia reports a version of it against contractualism, along with replies. The other side has its own hard cases, such as one person's agony for a hundred years against a million people's agony for a day less. The argument for the claim that violating ``utilitarianism'' leads to ``paradoxes, contradictions, circular preferences'' was in ``Circular Altruism'', which the book gives in shorter form as ``Feeling Moral''.

In short: a sound answer to Gowder on method, and good evidence that people fail to count lives, a point many of its opponents grant; on the disputed case, small harms against a grave one, the post depends on its argument against two tiers of value, a real objection that the other side has tried to answer.
\end{afterword}
